\chapter{The Extended and Error-State Kalman Filters}
\label{ch:eskf}

\begin{objectives}
\begin{itemize}
  \item Extend the Kalman filter to nonlinear models by linearization, and understand what that costs.
  \item Understand the three-way split into true, nominal and error state, and why it is the right way
        to filter on a manifold.
  \item Define measurement residuals and Jacobians on the manifold with the correct sign convention.
  \item Derive the injection and reset steps, including the reset Jacobian $G$.
  \item Assemble the complete error-state Kalman filter loop.
\end{itemize}
\end{objectives}

% ------------------------------------------------------------------------------
\section{The extended Kalman filter}

Real systems are nonlinear:
\begin{equation}
  x_k = f(x_{k-1}, u_k, w_k), \qquad z_k = h(x_k) + n_k,
\end{equation}
where $u_k$ is a known input such as an IMU reading. The \emph{extended} Kalman filter (EKF) propagates the
mean through the nonlinear functions and the covariance through their linearizations about the current
estimate:
\begin{align}
  &\hat x_k^- = f(\hat x_{k-1}^+, u_k, 0), &
  &P_k^- = F_k P_{k-1}^+ F_k\T + L_k Q_k L_k\T, \\
  &r_k = z_k - h(\hat x_k^-), &
  &H_k = \tfrac{\partial h}{\partial x}\big|_{\hat x_k^-},\quad K_k = P_k^- H_k\T(H_kP_k^-H_k\T + R_k)\inv,
\end{align}
with $F_k = \partial f/\partial x$ and $L_k = \partial f/\partial w$ evaluated at the estimate. The update then
follows \cref{thm:kf-update}.

The EKF is not optimal. It replaces each nonlinear transformation of a Gaussian by its linearization
(\cref{sec:linearization}), which is accurate only when the uncertainty is small compared with the
curvature of $f$ and $h$. Its Jacobians are evaluated at estimates that contain errors, which can make the
linearized model claim information the real system does not provide (the consistency problem of
\cref{ch:observability}). Despite this, the EKF and its variants dominate real-time navigation because they
are cheap and, when errors are kept small, accurate. The \emph{iterated} EKF relinearizes the measurement
function at the updated estimate and repeats the update, which is one Gauss--Newton step per iteration;
the unscented Kalman filter and particle filters avoid Jacobians altogether at a higher cost.

% ------------------------------------------------------------------------------
\section{Filtering on a manifold}

The EKF assumes $x$ lives in $\R^n$, so that $\hat x + K r$ makes sense. For VIO, $x$ contains a rotation.
Three quantities must therefore be distinguished.

\begin{definition}[True, nominal and error state]
\begin{itemize}
  \item The \textbf{true state} $x$ is the unknown physical state, a point on the state manifold
        $\mathcal{M}$.
  \item The \textbf{nominal state} $\hat x \in \mathcal{M}$ is the current best estimate, propagated with the
        full nonlinear model and no noise.
  \item The \textbf{error state} $\delta x = x \bminus \hat x \in \R^{n}$ is the small, minimal-dimension
        difference, defined with the operators of \cref{sec:boxplus}.
\end{itemize}
The filter's Gaussian belief is over the error state: $\delta x \sim \N(\delta\hat x, P)$, so that
$x = \hat x \bplus \delta x$.
\end{definition}

For the IMU state this means 16 nominal numbers (quaternion included) but a 15-dimensional error and a
$15\times15$ covariance, with no redundant directions and no constraints. The Kalman filter operates only on
the error state, which is a vector, so all of \cref{ch:kf} applies to it. This design is called the
\emph{error-state} (or \emph{indirect}) Kalman filter, ESKF for short. In aerospace the same idea for
attitude is the \emph{multiplicative} EKF.

\begin{keyidea}[why the error state]
\begin{itemize}
  \item \textbf{Minimal and unconstrained.} The error has exactly as many components as degrees of
        freedom, so its covariance is non-singular and needs no constraint handling.
  \item \textbf{Always near the origin.} After every update the error is reset to zero, so it never
        approaches the singularities of three-parameter rotations, and the small-angle approximations
        used to derive its dynamics remain valid.
  \item \textbf{Nearly linear.} The large, fast, nonlinear motion is carried by the nominal state, which is
        integrated exactly. The error evolves slowly and is small, so its dynamics are almost linear and
        linearization is accurate.
  \item \textbf{Slow.} Error dynamics are governed by sensor errors, not by vehicle motion, so the covariance
        could even be propagated at a lower rate than the nominal state if computation were scarce.
\end{itemize}
\end{keyidea}

% ------------------------------------------------------------------------------
\section{The ESKF cycle}

\begin{figure}[t]
\centering
\begin{tikzpicture}[node distance=14mm and 30mm]
  \node[block, minimum width=40mm] (nom) {\textbf{1. Nominal propagation}\\ $\hat x \gets f(\hat x, u)$, noise-free, nonlinear};
  \node[block, minimum width=40mm, right=of nom] (cov) {\textbf{2. Covariance propagation}\\ $P \gets \Phi P\Phi\T + Q_d$};
  \node[block, minimum width=40mm, below=of cov] (upd) {\textbf{3. Update}\\ $r = z\bminus h(\hat x)$, \ $\delta\hat x = Kr$\\ $P \gets$ Joseph form};
  \node[block, minimum width=40mm, below=of nom] (inj) {\textbf{4. Inject and reset}\\ $\hat x \gets \hat x\bplus\delta\hat x$, \ $\delta\hat x \gets 0$\\ $P \gets GPG\T$};
  \draw[->, thick] (nom) -- node[note, above]{every IMU} node[note, below]{sample} (cov);
  \draw[->, thick] (cov) -- node[note, right, align=left]{measurement\\arrives} (upd);
  \draw[->, thick] (upd) -- (inj);
  \draw[->, thick] (inj) -- (nom);
  \draw[->, thick, dashed] (cov.north) to[out=135, in=45] node[note, above]{no measurement: next IMU sample} (nom.north);\end{tikzpicture}
\caption{The error-state Kalman filter cycle. The error-state mean is zero throughout propagation, so only
the covariance is propagated in step 2. It becomes non-zero in step 3 and is transferred to the nominal
state and zeroed in step 4.}
\label{fig:eskf-cycle}
\end{figure}

\Cref{fig:eskf-cycle} shows the four steps.
\begin{enumerate}
  \item \textbf{Nominal propagation.} At each IMU sample, integrate the nonlinear kinematics with the
        measured inputs, ignoring noise (\cref{ch:strapdown}).
  \item \textbf{Error covariance propagation.} Propagate $P$ with the linearized error dynamics
        (\cref{ch:errordyn}). The error-state mean does not need propagating: it is zero after every reset,
        and the linear error dynamics driven by zero-mean noise keep it zero.
  \item \textbf{Update.} When a measurement arrives, compute the residual and its Jacobian on the manifold,
        and compute the Kalman estimate of the error, $\delta\hat x = Kr$, and the posterior covariance.
  \item \textbf{Injection and reset.} Move the estimated error into the nominal state with $\bplus$, set the
        error mean back to zero, and correct the covariance for the change of reference point.
\end{enumerate}

% ------------------------------------------------------------------------------
\section{Linearized error dynamics: the recipe}
\label{sec:error-dyn-recipe}

Suppose the true state obeys $\dot x = f(x, u, w)$ with noise $w$, and the nominal state obeys
$\dot{\hat x} = f(\hat x, u, 0)$. The error dynamics are derived in four mechanical steps.
\begin{enumerate}
  \item Write the true dynamics block by block.
  \item Substitute $x = \hat x\bplus\delta x$ in every block. For the attitude, $R = \hat R\Exp(\delta\theta)$.
  \item Subtract the nominal dynamics.
  \item Keep only terms of first order in the small quantities ($\delta x$ and $w$), using
        $\Exp(\delta\theta) \approx I + \sx{\delta\theta}$ and dropping their products.
\end{enumerate}
The result is a linear time-varying system
\begin{equation}
  \dot{\delta x} = F_c(\hat x, u)\,\delta x + G_c(\hat x)\,w,
\end{equation}
which \cref{sec:discretization} turns into $\Phi$ and $Q_d$. \Cref{ch:errordyn} applies the recipe to the
inertial navigation equations. For the attitude block, a short calculation is needed to express
$\frac{\dd}{\dd t}\Exp(\delta\theta)$, and \cref{ch:errordyn} works it out for the global perturbation as an
example and leaves the local one to you.

% ------------------------------------------------------------------------------
\section{Measurement updates on the manifold}
\label{sec:manifold-update}

Let a sensor produce a measurement $z$ on a manifold $\mathcal{Z}$ (a vector space, $\SO$, a pose, \dots)
according to the generative model
\begin{equation}
  z = h(x)\bplus n, \qquad n \sim \N(0, R_z).
\end{equation}
The ESKF compares the measurement with its prediction from the nominal state,
\begin{equation}
  r = z \bminus h(\hat x),
  \label{eq:manifold-residual}
\end{equation}
and models how the residual depends on the unknown true error $\delta x$. Substituting
$z = h(\hat x\bplus\delta x)\bplus n$ and linearizing,
\begin{equation}
  r \approx H\,\delta x + n, \qquad
  H = \frac{\partial}{\partial\delta x}\Big[ h(\hat x\bplus\delta x)\bminus h(\hat x) \Big]_{\delta x = 0}.
  \label{eq:manifold-H}
\end{equation}
$H$ is the manifold Jacobian of $h$ (\cref{def:manifold-jacobian}). With $r$ and $H$, the update is the
ordinary Kalman update of the error state with prior $\N(0, P)$:
\begin{equation}
  S = HPH\T + R_z,\qquad K = PH\T S\inv,\qquad \delta\hat x = Kr,\qquad P^+ = (I - KH)P(I - KH)\T + KR_zK\T.
\end{equation}

\begin{pitfall}[the sign of $H$]
There are two ways to set up the linearization, and they differ by a sign. The generative view above
differentiates the \emph{prediction} with respect to the \emph{true} error: $r \approx +H\delta x$. An
optimization view instead writes the residual as a function of a correction $\delta$ to the estimate,
$e(\delta) = z\bminus h(\hat x\bplus\delta)$, whose Jacobian is $-H$ (to first order). Both lead to the same
correction if used consistently. Mixing them flips every correction. In this project, $H$ is always
\eqref{eq:manifold-H}, and its finite-difference test perturbs the nominal state:
$r(\delta x) := h(\hat x\bplus\delta x)\bminus h(\hat x)$ and $H \approx \partial r/\partial\delta x$.
\end{pitfall}

For a vector-valued measurement, $\bminus$ is subtraction and this reduces to the EKF. For a
rotation-valued measurement such as the pose measurement of Milestone 6, the attitude part of the residual
is $\Log_q(\hat q\inv\qprod q_{meas})$ (AGENTS.md F3).

% ------------------------------------------------------------------------------
\section{Injection and reset}
\label{sec:reset}

After the update, the error estimate $\delta\hat x = Kr$ is non-zero. Injection moves it into the nominal
state:
\begin{equation}
  \hat x^+ = \hat x\bplus\delta\hat x:\quad
  \hat p \gets \hat p + \delta\hat p,\quad \hat v \gets \hat v + \delta\hat v,\quad
  \hat q \gets \hat q\qprod\Exp_q(\delta\hat\theta),\quad \hat b \gets \hat b + \delta\hat b,
\end{equation}
followed by normalization of every quaternion that was touched, for the IMU and every clone (AGENTS.md F4).

The error is now defined relative to the new nominal state, $\delta x^+ = x\bminus\hat x^+$, and its mean is zero.
But its covariance is not exactly $P^+$, because the reference point of the tangent space has moved.

\begin{proposition}[Reset Jacobian]
\label{prop:reset}
Let $\delta x \sim \N(\delta\hat x, P^+)$ be the error relative to $\hat x$ after the update, and
$\delta x^+ = x\bminus(\hat x\bplus\delta\hat x)$ the error relative to the injected nominal state. To first order,
$\delta x^+ \sim \N(0, GP^+G\T)$ with $G = \partial\delta x^+/\partial\delta x$ evaluated at $\delta\hat x$. For the vector
blocks $G = I$, and for each attitude block
\begin{equation}
  G_\theta = I - \tfrac12\sx{\delta\hat\theta}.
  \label{eq:reset-G}
\end{equation}
\end{proposition}
\begin{proof}
For vector blocks $\delta x^+ = \delta x - \delta\hat x$, whose Jacobian is $I$. For the attitude, $q = \hat q\qprod\Exp_q(\delta\theta)$
and $q = \hat q^+\qprod\Exp_q(\delta\theta^+)$ with $\hat q^+ = \hat q\qprod\Exp_q(\delta\hat\theta)$, so
\begin{equation*}
  \Exp(\delta\theta^+) = \Exp(-\delta\hat\theta)\,\Exp(\delta\theta).
\end{equation*}
By the BCH formula \eqref{eq:bch}, $\delta\theta^+ = \delta\theta - \delta\hat\theta - \tfrac12\,\delta\hat\theta\times\delta\theta + \cdots$,
whose derivative with respect to $\delta\theta$ is $I - \tfrac12\sx{\delta\hat\theta}$. At $\delta\theta = \delta\hat\theta$ the value is
exactly zero, so the new error has zero mean.
\end{proof}

\begin{remark}[relation to the right Jacobian]
The derivative of $\Log(\Exp(-\delta\hat\theta)\Exp(\delta\theta))$ at $\delta\theta = \delta\hat\theta$ can be computed exactly with
\eqref{eq:jr-def}: it is $\Jr(\delta\hat\theta) = I - \tfrac12\sx{\delta\hat\theta} + O(\|\delta\hat\theta\|^2)$. The frozen
specification (AGENTS.md F4) uses the first-order form \eqref{eq:reset-G}. For a typical correction of a few
milliradians the two differ by about $10^{-6}$ and the difference is irrelevant. Changing it would be a
logged convention change, not a silent substitution.
\end{remark}

The covariance reset is therefore
\begin{equation}
  P \gets G P G\T, \qquad G = \blkdiag\big(I_6,\ I - \tfrac12\sx{\delta\hat\theta},\ I_6,\ \ldots\big),
\end{equation}
with one $I - \tfrac12\sx{\delta\hat\theta_i}$ block for each clone's attitude and identity elsewhere.

\begin{convention}[injection and reset, AGENTS.md F4]
Inject with $\bplus$ into the IMU block and every clone; normalize every touched quaternion; reset
$\delta\hat x \gets 0$ and $P \gets GPG\T$ with $G$ implemented exactly as above. Using $G \approx I$ is a deviation
that must be logged.
\end{convention}

How much does $G$ matter? It rotates the attitude covariance by half the injected correction, a
first-order effect in $\|\delta\hat\theta\|$. With small corrections at high rate the effect is tiny, and many
implementations use $G = I$. Large corrections (after a long outage, or a loop closure) make it
noticeable. It costs almost nothing, and implementing it exactly removes one approximation from the
consistency analysis of Milestone 7.

% ------------------------------------------------------------------------------
\section{The complete error-state filter}

\begin{algorithm}[t]
\caption{Error-state Kalman filter (one IMU step and an optional update)}
\label{alg:eskf}
\begin{algorithmic}[1]
\Require nominal $\hat x$, covariance $P$, IMU sample at $t_k$, optional measurement $z$ at $t_k$
\State $\Dt \gets t_k - t_{k-1}$ \Comment{from integer nanosecond timestamps}
\State $\hat x \gets f(\hat x, u_k, \Dt)$; normalize $\hat q$ \Comment{nominal propagation}
\State $\Phi \gets I + F_c\Dt$; \quad $Q_d \gets G_cQ_cG_c\T\Dt$
\State $P \gets \Phi P\Phi\T + Q_d$; \quad $P \gets \tfrac12(P + P\T)$
\If{measurement $z$ available}
  \State $r \gets z\bminus h(\hat x)$; \quad $H \gets \partial[h(\hat x\bplus\delta x)\bminus h(\hat x)]/\partial\delta x$
  \State $S \gets HPH\T + R_z$; factor $S$; gate on $r\T S\inv r$
  \State $K \gets PH\T S\inv$ (by solving); \quad $\delta\hat x \gets Kr$
  \State $P \gets (I-KH)P(I-KH)\T + KR_zK\T$
  \State $\hat x \gets \hat x\bplus\delta\hat x$; normalize quaternions \Comment{inject}
  \State $P \gets GPG\T$; \quad $P \gets \tfrac12(P + P\T)$ \Comment{reset}
\EndIf
\end{algorithmic}
\end{algorithm}

\Cref{alg:eskf} assembles the pieces. It is generic: \cref{ch:strapdown,ch:errordyn} supply $f$, $F_c$ and $G_c$
for the IMU, \cref{ch:updates} supplies $h$ and $H$ for a pose sensor, and \cref{ch:msckfupdate} supplies them for
a camera. When measurements arrive between IMU samples, the filter first propagates to the measurement time
by splitting the IMU interval (\cref{ch:strapdown}).

\begin{projectlink}[Milestones 4--6]
M4 implements step 2 of \cref{alg:eskf}, M5 steps 3--4, and M6 the update, injection and reset with a pose
measurement. Testing them separately localizes bugs: if propagation is exact on synthetic data, and the
covariance matches Monte Carlo statistics, then any later inconsistency is in the update.
\end{projectlink}

\begin{takeaways}
\begin{itemize}
  \item The EKF linearizes about the estimate; it is accurate when errors are small relative to curvature.
  \item The ESKF splits the state into a nonlinear nominal part and a small, minimal, nearly linear error
        part. The Kalman filter only ever sees the error.
  \item The error mean is zero during propagation; only $P$ is propagated.
  \item Residual $r = z\bminus h(\hat x) \approx H\delta x + n$ with $H$ the manifold Jacobian of $h$ at $\hat x$.
  \item After injection, reset $P \gets GPG\T$ with $G_\theta = I - \tfrac12\sx{\delta\hat\theta}$ for every attitude block.
\end{itemize}
\end{takeaways}

\begin{exercises}
  \item Write the EKF for a planar robot with state $(x, y, \psi)$ driven by speed and turn rate and measured by
        range to a known beacon. Identify where linearization error enters.
  \item Show that if the error state is zero-mean after a reset and $\dot{\delta x} = F_c\delta x + G_cw$ with
        zero-mean $w$, then it stays zero-mean until the next update.
  \item Derive the reset Jacobian for the \emph{global} perturbation $R = \Exp(\delta\theta)\hat R$. How does it differ
        from \eqref{eq:reset-G}?
  \item Suppose $\delta\hat\theta = (0, 0, 0.1)$~rad and $P_{\theta\theta} = \diag(10^{-4}, 10^{-4}, 10^{-2})$. Compute
        $GP_{\theta\theta}G\T$ and compare with $\Jr(\delta\hat\theta)P_{\theta\theta}\Jr(\delta\hat\theta)\T$. When is the difference
        noticeable?
  \item An iterated ESKF repeats the update with $H$ re-evaluated at $\hat x\bplus\delta\hat x$. Show that one
        iteration is a Gauss--Newton step on the MAP cost $\|\delta x\|^2_{P\inv} + \|z\bminus h(\hat x\bplus\delta x)\|^2_{R_z\inv}$.
\end{exercises}
